\section*{Chapitre \ref{ch:dv:american-options-and-early-exercise} --- Options américaines et exercice anticipé}

\begin{solution}{exo:dv:american-options-and-early-exercise:1}
$6.0902-5.5735=0.5167$, soit 9.3~\% de la valeur européenne.
\end{solution}

\begin{solution}{exo:dv:american-options-and-early-exercise:2}
$\gamma=2\times0.05/0.04=2.5$, $S^*=100\times2.5/3.5=71.43$~; à 100, le put vaut
$28.57\times(100/71.43)^{-2.5}=12.32$.
\end{solution}

\begin{solution}{exo:dv:american-options-and-early-exercise:3}
En deçà de l'intérêt couru sur le strike pendant les 60 jours, $50(1-e^{-0.05\times60/365})=0.41$~:
l'exercice fait alors perdre au moins l'intérêt.
\end{solution}

\begin{solution}{exo:dv:american-options-and-early-exercise:4}
L'exercice anéantit la valeur résiduelle de l'option, qui croît avec la volatilité, et rapporte
l'intérêt sur le strike, qui croît avec le taux. Des frais d'emprunt agissent comme un
taux de dividende~: ils abaissent le forward, accroissent la valeur d'attendre avec un put et rendent ainsi
son exercice anticipé moins probable~; sur un call, ils peuvent le rendre intéressant.
\end{solution}

\begin{solution}{exo:dv:american-options-and-early-exercise:5}
Avec une seule date, à l'échéance, le porteur n'a pas le choix~: c'est l'option européenne. Ajouter
une date ajoute un choix~: le porteur peut toujours l'ignorer, si bien que la valeur ne peut pas baisser
(le nouveau supremum porte sur un ensemble plus grand de temps d'arrêt).
\end{solution}

\begin{solution}{exo:dv:american-options-and-early-exercise:6}
Le put profondément dans la monnaie à 8~\%, d'erreur $-0.064$~: là, la \omterm{def:dv:american-options-and-early-exercise:premium}{prime d'exercice anticipé} est grande et
l'option est exercée rapidement, alors que la loi de puissance unique de l'approximation est calée sur une
forme perpétuelle et mal évalue la \omterm{def:dv:american-options-and-early-exercise:boundary}{frontière d'exercice} proche.
\end{solution}

\begin{solution}{exo:dv:american-options-and-early-exercise:7}
0.932, contre un intérêt sur le strike de 0.785~: avec 90 jours restants, la composante put du call
vaut 0.15, et le dividende doit aussi la compenser.
\end{solution}

\begin{solution}{exo:dv:american-options-and-early-exercise:8}
La formule européenne ne comporte pas de \omterm{def:dv:american-options-and-early-exercise:premium}{prime d'exercice anticipé}, si bien qu'il lui faut une volatilité plus élevée
(21.4~\%) pour atteindre le prix américain. Évalué avec un modèle américain, le put se situe à 20~\%,
comme les options européennes. La «~richesse~» est la prime, et non la volatilité.
\end{solution}

\begin{solution}{pb:dv:american-options-and-early-exercise:1}
\textbf{1.} 20.30.
\textbf{2.} 19.56.
\textbf{3.} Oui~: 0.74 de mieux par action.
\textbf{4.} $80(1-e^{-0.04\times29/365})=0.254$.
\textbf{5.} 0.255~: le call est si profondément dans la monnaie que sa valeur temps est l'intérêt sur
le strike augmenté d'un put valant un dixième de cent.
\textbf{6.} $0.7446\times100\times10\,000$~: 744\,600~USD.
\textbf{7.} 45~\% du total~: 335\,000~USD.
\textbf{8.} Les vendeurs non assignés~: la chambre de compensation assigne les notifications d'exercice
au hasard aux membres compensateurs, qui les assignent à leurs clients~; un vendeur
resté non assigné conserve un call valant 19.56 qu'il a vendu sur une action valant 20.30.
\textbf{9.} Elle ne sait pas combien de porteurs exerceront, ni quels vendeurs seront
assignés~: le transfert est une espérance, et non un flux couvrable.
\textbf{10.} Elle doit livrer des actions à 80~; si elle détenait des actions en couverture, elle les livre
et reçoit 80, mais ne perçoit pas le dividende afférent.
\textbf{11.} Égal ou légèrement inférieur à 20.30~: nul ne devrait payer plus que la valeur d'exercice, et
les porteurs qui peuvent exercer ne devraient pas vendre en dessous.
\textbf{12.} Environ 19.56.
\textbf{13.} Investis $100.30-20.30=80$~; après la date de détachement, elle détient le dividende 1.00, une
action à 99.30 et un call vendu valant 19.56~: $1.00+99.30-19.56-80=+0.74$.
\textbf{14.} Zéro, hors frais~: l'action est livrée à 80, le montant investi.
\textbf{15.} Par le taux d'échec attendu et son incertitude, sachant que le gain n'est versé
que sur les calls non assignés, et que l'ensemble du marché joue la même opération.
\textbf{16.} L'inattention, la détention via un courtier sans instructions, et la croyance que
les options sont exercées automatiquement avant l'échéance.
\textbf{17.} Exercer pour le compte du client exige son accord et son financement (le
client doit payer le strike)~; les courtiers avertissent plutôt que d'exercer.
\textbf{18.} Comme si chaque porteur exerçait de manière optimale (la valeur prudente), ce qui libère le gain
sur les calls non assignés.
\textbf{19.} Un dividende-seuil de 0.255~; 335\,000~USD transférés aux vendeurs.
\textbf{20.} Lorsqu'un dividende versé avant l'échéance vaut davantage que l'intérêt sur le strike
pendant la durée de vie restante, augmenté de la composante put du call (ou lorsque des frais
d'emprunt rendent précieuse la détention de l'action).
\end{solution}

\begin{solution}{iq:dv:american-options-and-early-exercise:1}
Un call~: seulement juste avant une date de détachement, lorsque le dividende dépasse l'intérêt sur le
strike augmenté de l'optionalité perdue (ou sur une action difficile à emprunter). Un put~: oui, lorsqu'il est profondément dans la
monnaie avec des taux positifs, puisque l'exercice rapporte l'intérêt sur le strike.
\iqlookfor{l'argument de l'intérêt sur le strike dans les deux sens.}
\end{solution}

\begin{solution}{iq:dv:american-options-and-early-exercise:2}
Le porteur choisit quand arrêter, en ne s'appuyant que sur l'information disponible à l'instant, et
choisira la meilleure règle~: la valeur est le supremum sur les temps d'arrêt de l'espérance actualisée
du gain sous la mesure de valorisation. Un arbre le calcule à rebours, en prenant à chaque nœud
le maximum entre la valeur d'exercice et l'espérance actualisée de l'étape suivante
(l'enveloppe de Snell, livre~4, chapitre~10).
\iqlookfor{arrêt optimal, et récurrence rétrograde comme algorithme.}
\end{solution}

\begin{solution}{iq:dv:american-options-and-early-exercise:3}
Équation de valorisation indépendante du temps $\tfrac12\sigma^2S^2V^{\prime\prime}+rSV^{\prime}-rV=0$~; solutions $S$ et
$S^{-\gamma}$ avec $\gamma=2r/\sigma^2$~; le put conserve $AS^{-\gamma}$~; raccord de valeur et
recollement C1 (smooth pasting) en $S^*$ donnent $S^*=K\gamma/(1+\gamma)$.
\iqlookfor{les deux conditions aux limites.}
\end{solution}

\begin{solution}{iq:dv:american-options-and-early-exercise:4}
Les porteurs devraient exercer ce soir~; certains ne le feront pas. Prévoir des assignations sur la plupart, mais non sur la totalité de la
position~; la couverture en actions doit être dimensionnée pour la livraison (détenir les actions), tandis que le dividende
sur les actions ainsi livrées est perdu~; les calls non assignés valent moins que leur
valeur d'exercice le lendemain, un gain sur lequel il ne faut pas compter.
\iqlookfor{incertitude sur les assignations et son effet sur la couverture.}
\end{solution}

\begin{solution}{iq:dv:american-options-and-early-exercise:5}
Une approximation fermée (quadratique ou Bjerksund--Stensland) ou une grille précalculée dans
la boucle de cotation~; un arbre ou une grille aux différences finies, avec dividendes, pour le risque et la fin de
journée. Vérifier l'approximation face au pricer exact sur toute la gamme de moneyness,
les échéances et les taux, en particulier profondément dans la monnaie et près des dates de détachement, et se rabattre sur la
grille là où elle échoue.
\iqlookfor{rapidité là où elle est nécessaire, précision là où elle compte, et un indice de référence.}
\end{solution}

\begin{solution}{iq:dv:american-options-and-early-exercise:6}
Inverser le pricer américain (arbre ou grille) avec les dividendes et l'emprunt, ou retrancher une
estimation de la \omterm{def:dv:american-options-and-early-exercise:premium}{prime d'exercice anticipé} et inverser la formule européenne. Avec la
formule européenne, la prime est lue comme une volatilité~: les puts et calls profondément dans la
monnaie, avant dividendes, affichent des volatilités trop élevées, un faux smile de volatilité.
\iqlookfor{la prime prise pour de la volatilité.}
\end{solution}
